\section{Succinct Representation (SODA 2015, §3.2)}
\label{sec:ere}

The inner layer of every approximate filter is an \emph{Exact Range Emptiness} (ERE)
structure: given a universe $[r]$ of $\ell = \lceil \log_2 r \rceil$-bit keys
and a set $S' \subseteq [r]$ of $n$ keys, it answers
``$[a, b] \cap S' = \emptyset$?'' exactly (no false positives or false negatives).
Throughout this chapter, $S'$ denotes the input to ERE; in the ARE pipeline
(Chapter~\ref{chap:hash}), $S' = h(S)$.

\subsection{Two-Level Hierarchy}

To achieve the information-theoretic minimum of $n \log_2(r/n) + O(n)$ bits, the
structure divides the universe $[r]$ into $n$ equal-sized blocks:

\paragraph{Global level (succinct indexing).}
Two bit-vectors provide $O(1)$ navigation:
\begin{itemize}
  \item $D_1$ (size $n$ bits): $D_1[i] = 1$ if block $i$ contains at least one key.
  \item $D_2$ (size ${\sim}2n$ bits): an Elias--Fano~\citep{elias1974efficient,fano1971number}
    representation of block counts. Each non-empty block $i$ contributes a $1$ followed
    by $n_i$ zeros. A \textsc{Select}$_1$ / \textsc{Rank}$_0$ pair maps any block index
    to its starting position in the data array in $O(1)$ time.
\end{itemize}

\paragraph{Local level (bit-packed suffixes).}
Instead of storing full keys, only the suffix $x \bmod (r/n)$ is stored for each
key~$x$. The suffix length is $W = |r| - \lceil \log_2 n \rceil$ bits. All suffixes are packed into a dense \texttt{uint64} array using bit-shifts,
eliminating per-element overhead.

\subsection{Local Search in Buckets}
\label{sec:ere-local-search}

The original paper~\citep{goswami2015approximate} suggests a Weak Prefix Search structure
(Hollow Z-Fast Trie) for $O(1)$ worst-case local queries.
We use binary search on bit-packed suffixes instead, with an adaptive
fallback to linear scan for small buckets.
Section~\ref{sec:ere-no-wps} details why.

\paragraph{Block occupancy under uniform distribution.}
The structure divides the universe into $m = 2^{\lfloor \log_2 n \rfloor}$ blocks;
each non-empty block stores its keys' suffixes in a packed array (a \emph{bucket}).
The expected
number of keys per block is $\lambda = n/m \in [1,\,2)$.
When $n$ is a power of two, $m = n$ and $\lambda = 1$ exactly.
Each key falls into a uniformly random block, so the occupancy of a fixed block
follows $\mathrm{Bin}(n,\,1/m)$, which converges to $\mathrm{Poisson}(\lambda)$
as $n \to \infty$~\citep[Ch.~5]{mitzenmacher2017probability}.

For $\lambda = 1$, the fraction of empty blocks converges to $e^{-1} \approx 36.8\%$,
and the mean occupancy among non-empty blocks to $1/(1 - e^{-1}) \approx 1.58$.

To estimate the maximum bucket size, we use a union bound.
The probability that a fixed block contains at least $k$ keys is
$\Pr[\mathrm{Poisson}(\lambda) \geq k] = 1 - \sum_{j=0}^{k-1} e^{-\lambda}\lambda^j/j!$
(which simplifies to $1 - \sum_{j=0}^{k-1} e^{-1}/j!$ when $\lambda = 1$).
The probability that \emph{any} of the $m$ blocks contains $\geq k$ keys
is at most
\[
  \Pr\!\bigl[\max_i X_i \geq k\bigr]
    \;\leq\; m \cdot \Pr[\mathrm{Poisson}(\lambda) \geq k].
\]
The expected number of blocks with occupancy $\geq k$ drops below~1 when
$m \cdot \Pr[\mathrm{Poisson}(\lambda) \geq k] < 1$.
In the tables below we use $\lambda = 1$ ($m = n$, powers of two).
Table~\ref{tab:ere-bucket-stats} shows that this threshold closely matches
the measured maximum.

\begin{table}[h]
\centering
\begin{tabular}{@{}rrrrrr@{}}
\toprule
$n$ & Empty (\%) & Avg / non-empty & Max (meas.) &
  $k^*$ & $n \cdot \Pr[X \geq k^*]$ \\
\midrule
$2^{20}$ & 36.75 & 1.58 & 9  & 10 & 0.12 \\
$2^{24}$ & 36.79 & 1.58 & 9  & 11 & 0.17 \\
$2^{27}$ & 36.79 & 1.58 & 10 & 12 & 0.11 \\
$2^{30}$ & 36.79 & 1.58 & 13 & 12 & 0.89 \\
\bottomrule
\end{tabular}
\caption{ERE bucket statistics for uniform keys with $\lambda = 1$ ($m = n$).
  $k^*$ is the smallest $k$ such that $n \cdot \Pr[\mathrm{Poisson}(1) \geq k] < 1$.
  The measured maximum can exceed $k^*$ (e.g.\ $n = 2^{30}$: max $= 13 > k^* = 12$)
  because the bound is not tight;
  $n \cdot \Pr[\mathrm{Poisson}(1) \geq 13] \approx 0.07$ at $n = 2^{30}$
  (see Table~\ref{tab:ere-poisson-tail}).}
\label{tab:ere-bucket-stats}
\end{table}

Table~\ref{tab:ere-poisson-tail} shows how rapidly the tail probability
drops with $k$: even at $n = 2^{30}$, the probability of any block containing
$\geq 20$ keys is below $10^{-10}$, making buckets larger than ${\sim}15$
keys virtually impossible for uniform data.

\begin{table}[h]
\centering
\begin{tabular}{@{}rrr@{}}
\toprule
$k$ & $\Pr[\mathrm{Poisson}(1) \geq k]$ &
  $n \cdot \Pr[X \geq k]$,\; $n = 2^{30}$ \\
\midrule
10 & $1.1 \times 10^{-7}$  & $1.2 \times 10^{2\phantom{0}}$ \\
13 & $6.4 \times 10^{-11}$ & $6.8 \times 10^{-2}$ \\
15 & $3.0 \times 10^{-13}$ & $3.2 \times 10^{-4}$ \\
20 & $1.6 \times 10^{-19}$ & $1.7 \times 10^{-10}$ \\
25 & $2.5 \times 10^{-26}$ & $2.6 \times 10^{-17}$ \\
\bottomrule
\end{tabular}
\caption{Poisson tail probabilities for bucket occupancy ($\lambda = 1$),
  computed via $\Pr[X \geq k] = 1 - \sum_{j=0}^{k-1} e^{-1}/j!$.
  At $k = 13$ the expected number of such blocks across all $n = 2^{30}$
  blocks is ${\sim}0.07$ --- consistent with observing max $= 13$ occasionally.
  At $k = 20$ it is below $10^{-10}$ even for $n = 2^{30}$.}
\label{tab:ere-poisson-tail}
\end{table}

\paragraph{Non-uniform distributions and worst-case bounds.}
For non-uniform input, the Poisson model does not apply: clustered keys
concentrate in a few blocks. When ERE is used inside SODA~ARE
(Section~\ref{sec:soda}), the locality-preserving hash $h(x) = (u + x) \bmod 2^K$ (where $u$ is a
per-block constant) maps consecutive keys to consecutive hashed values,
preserving local density.

However, the bucket size is always bounded by the block capacity $2^w$,
where $w = K - \lfloor \log_2 n \rfloor$ is the suffix length
and $K$ is the fingerprint width in bits (defined in Section~\ref{sec:soda}).
For industry-standard space budgets (BPK~$\leq 15$; since BPK~$\approx w + 3.2$
from Table~\ref{tab:ere-space}, this gives $w \leq 12$),
so the maximum bucket contains at most $2^{12} = 4096$ keys.
At $w = 12$ bits, such a worst-case bucket occupies
$4096 \times 12 / 8 = 6144$~bytes --- well within L1 cache (typically 32--64~KB).

\paragraph{Search strategy.}
We implement both binary search ($O(\log k)$) and linear scan ($O(k)$) over
bit-packed suffixes in a bucket of $k$ keys.
Benchmarks on Apple~M4~Max show that linear scan is ${\sim}3\times$ faster
for $k \leq 12$ (1.3--2.7~ns vs.\ 5--9~ns) due to sequential prefetch,
while binary search wins for $k > 128$ (21~ns vs.\ 35~ns at $k = 256$).
The crossover is at $k \approx 128$ and is stable across suffix widths
$w \in \{8, 11, 16, 20\}$.

The default query uses adaptive dispatch: linear scan for buckets with
$k \leq 128$, binary search above.
The total ERE query takes ${\sim}36$~ns on Apple~M4~Max (stable across
$n \in [2^{10},\, 2^{20}]$ and key widths 64--512 bits), of which
${\sim}27$--$35$~ns is Rank/Select overhead on $D_1$ and $D_2$.
Bucket search contributes at most ${\sim}9$~ns (binary, $k = 12$) or
${\sim}2.7$~ns (linear, $k = 12$) --- roughly 7--25\% of the total.
This makes the choice between search strategies a micro-optimization;
the key property is that even in the worst case ($k = 2^w$, binary search),
query time grows only logarithmically.

\subsection{Why Not Weak Prefix Search}
\label{sec:ere-no-wps}

\citet{goswami2015approximate} propose $O(1)$ worst-case local queries via a Weak Prefix
Search structure built from a Hollow Z-Fast Trie (HZFT) and Monotone Minimal
Perfect Hashing (MMPH). The authors analyse these components in a word-RAM model
where hash table lookups and trie navigation cost $O(1)$, and the space is
expressed up to $O(n)$ additive terms. This is sufficient for a theoretical result,
but the hidden constants become dominant in practice.

We implemented three variants of Weak Prefix Search, collectively referred to as
LERLOC (\textbf{L}ocal \textbf{E}xact \textbf{R}ange \textbf{LOC}ator),
differing in the trie and MMPH components used:

\begin{table}[h]
\centering
\begin{tabular}{@{}lrr@{}}
\toprule
Variant & Space (bits/key) & Query (ns), $L{=}64$ \\
\midrule
LERLOC Fast (HZFT + BoomPHF)     & 179 & 1095--1466 \\
LERLOC Compact (SHZFT + BoomPHF) & 151 & 372--567   \\
LeMon-LERLOC (SHZFT + LeMonHash) &  59 & 537--762   \\
\bottomrule
\end{tabular}
\caption{LERLOC variants: space and query latency on Apple~M4~Max,
  $n = 1024$--$32768$, 64-bit keys.
  LeMon-LERLOC breakdown: SHZFT trie ${\sim}35$~BPK, LeMonHash MMPH ${\sim}19$~BPK,
  RSDic leaf bitvector ${\sim}5$~BPK.
  The MMPH indexes the boundary set $P$ of the trie: 3 strings per node
  ($n$ leaves $+ (n{-}1)$ internal), so $|P| \leq 3(2n - 1) = 6n - 3$
  before deduplication; empirically $|P|/n \approx 4.3$ for $n > 10^5$.}
\label{tab:lerloc-cost}
\end{table}

The high space cost stems from three expansion factors hidden behind
$O(n)$ in the paper:
(1)~the SHZFT hash table stores not only $2n{-}1$ true descriptors but also
${\sim}3.5n$ pseudo-descriptors required by Fat Binary Search, totalling
${\sim}5.5n$ entries at ${\sim}35$~BPK;
(2)~the range locator boundary set $P$ contains 3~strings per trie node
($\leq 6n{-}3$ before deduplication, empirically ${\sim}4.3n$);
(3)~the MMPH on $P$ stores a \texttt{uint8} rank per element ($8$~bits
$\times\,4.3 \approx 34$~BPK).
These factors stack: the total ``$O(n)$ additional space'' becomes
${\sim}100$~bits/key in practice (see the detailed breakdown in
\texttt{locators/rloc/MEMORY\_INVESTIGATION.md}).

When ERE is used inside an ARE filter,
the suffix width is $w = K - \lfloor \log_2 n \rfloor \approx \lceil \log_2(\mathcal{L}/\varepsilon) \rceil$
(since $K = \lceil \log_2(n\mathcal{L}/\varepsilon) \rceil$),
which depends only on the range length $\mathcal{L}$ and FPR $\varepsilon$,
not on $n$. For typical parameters ($\mathcal{L} \leq 64$,
$\varepsilon \geq 0.01$), $w$ is $7$--$12$ bits.
The data itself costs $w$ bits per key; even the most compact LERLOC variant (59~BPK)
adds $5$--$8 \times$ more space than the data it indexes.

Query latency is equally problematic. LERLOC replaces only the local search
inside a bucket, while Rank/Select navigation on $D_1$ and $D_2$
(${\sim}27$~ns) remains unchanged.
Each LERLOC call descends the trie (pointer chasing through hash tables), then
performs an MMPH rank lookup --- at least 370~ns for the Compact variant at $L = 64$.
Binary search on a typical bucket of ${\sim}12$ keys takes 5--9~ns
(Section~\ref{sec:ere-local-search}), making LERLOC $40$--$110 \times$
slower for the operation it replaces.

The root cause is that the $O(1)$ theoretical guarantee hides large constants:
each ``constant-time'' trie step involves hashing a variable-length prefix,
probing a hash table, and following a pointer --- operations that are
individually cheap but accumulate across $O(\log \log r)$ trie levels.
For the small buckets that arise in practice (${\leq}13$ keys for uniform
data, ${\leq}2^w$ in general), binary search on packed suffixes is both
faster and requires no additional space.

Even in an adversarial worst case --- every key in a single bucket ---
binary search remains fast:

\begin{table}[h]
\centering
\begin{tabular}{@{}rrr@{}}
\toprule
$w$ (suffix bits) & Bucket size ($2^w$) & Binary search (ns) \\
\midrule
12 & 4\,096   & 37 \\
15 & 32\,768  & 48 \\
17 & 131\,072 & 58 \\
\bottomrule
\end{tabular}
\caption{Worst-case binary search on a fully packed bucket (Apple~M4~Max).
  At $w = 12$ (BPK~$\approx 15$), the entire bucket search costs 37~ns ---
  still $10 \times$ faster than the best LERLOC variant.}
\label{tab:ere-worst-case-binsearch}
\end{table}

\subsection{Space Analysis}

\begin{table}[h]
\centering
\begin{tabular}{@{}lrrr@{}}
\toprule
$\ell$ (key bits) & Suffix bits ($|r| - \log_2 n$) & Metadata & Total BPK \\
\midrule
64   & 44.07  & 3.2 & 47.27  \\
128  & 108.07 & 3.2 & 111.27 \\
256  & 236.07 & 3.2 & 239.27 \\
512  & 492.07 & 3.2 & 495.27 \\
\bottomrule
\end{tabular}
\caption{Measured ERE space with $n = 10^6$ keys. Metadata ($D_1 + D_2$ + rank/select
  indices) is constant at ${\sim}3.2$ BPK.}
\label{tab:ere-space}
\end{table}

Space grows linearly with $\ell$ --- unavoidable for an exact structure, since it must
store the information distinguishing the keys. This linear dependence motivates the
move to approximate range emptiness: by storing hashed fingerprints of length
$K \approx \log_2(\mathcal{L}/\varepsilon)$ instead of full keys, the space becomes
independent of the original key length~$\ell$.

\subsection{Complexity Summary}

\begin{itemize}
  \item \textbf{Build:} $O(n \cdot |r|)$ --- a single pass over sorted keys.
  \item \textbf{Query:} $O(|r|/w)$ expected (block index + binary search over ${\sim}2$
    keys of $\ell$ bits each); $O((|r|/w) \log n)$ worst case if all keys collide in a
    single block. When used inside ARE with $K$-bit fingerprints fitting in one machine
    word, this reduces to $O(1)$ expected.
  \item \textbf{Space:} $n \log_2(r/n) + O(n)$ bits, matching the information-theoretic
    lower bound.
\end{itemize}
